\section{KernelSHAP: hồi quy để tính giá trị Shapley}
\label{sec:ch04-kernelshap}

Sau khi đã chọn cách gán giá trị cho liên minh, ta vẫn phải tính các giá trị
\tn{giá trị Shapley}{Shapley value} đó. Khi có $M$ đặc trưng, công thức~\ref{eq:ch04-gia-tri-shapley} đòi hỏi xem nhiều
liên minh. KernelSHAP viết cùng bài toán dưới dạng hồi quy có trọng số trên các
vector $z'$ biểu thị liên minh. Mô hình hồi quy vẫn là
phương trình~\ref{eq:ch04-mo-hinh-cong-tinh}, nhưng trọng số của một liên minh
phụ thuộc số đặc trưng nó chứa. Các liên minh rất nhỏ và rất lớn được nhấn
mạnh để ràng buộc giá trị nền và tổng các đóng góp.

Hình~\ref{fig:ch04-kernelshap} cho thấy phần hồi quy chỉ là lớp tính toán. Mỗi
hàng đầu vào của nó trước hết phải được đánh giá bằng quy tắc hoàn thiện đặc
trưng ở mục~\ref{sec:ch04-gia-tri-lien-minh-va-phan-phoi-nen}. Vì vậy,
KernelSHAP không loại bỏ giả định về dữ liệu; nó dùng một \tn{hàm nhân Shapley}{Shapley
kernel} để làm hồi quy trả về cách chia mà các tiên đề đã chọn.

\begin{figure}[htbp]
  \centering
  \begin{tikzpicture}[
  box/.style={draw=black!65, rounded corners=2pt, align=center,
    minimum width=4.0cm, minimum height=0.9cm, fill=black!5},
  data/.style={draw=black!65, rounded corners=2pt, align=center,
    minimum width=4.0cm, minimum height=0.9cm, fill=black!15},
  arrow/.style={-{Stealth[length=2mm]}, draw=black!65, thick}
]
  \node[data] (coalitions) {Liên minh $z'$ của các đặc trưng};
  \node[box, below=8mm of coalitions] (value) {Hoàn thiện đặc trưng thiếu\\
và tính $v(S)$};
  \node[data, right=15mm of value] (kernel) {Hàm nhân Shapley\\
theo $|S|$};
  \node[box, below=10mm of value, xshift=28mm] (regression)
    {Hồi quy có trọng số\\
để tìm $\phi_0,\ldots,\phi_M$};

  \draw[arrow] (coalitions) -- (value);
  \draw[arrow] (value) -- (regression);
  \draw[arrow] (kernel.south) |- (regression.east);
\end{tikzpicture}

  \caption{KernelSHAP đi từ các liên minh đặc trưng tới một hồi quy có trọng số.
  Giá trị của mỗi liên minh phụ thuộc quy tắc hoàn thiện phần đặc trưng còn
  thiếu; kernel Shapley chỉ quyết định cách các liên minh ấy ràng buộc các hệ
  số $\phi_i$. Cách viết hồi quy này bắt nguồn từ Lundberg và Lee
  ~\autocite{p10shap}.}
  \label{fig:ch04-kernelshap}
\end{figure}

Điểm này nối thẳng với LIME. Cả hai đều có mô hình cộng tính và có thể dùng
hồi quy bình phương tối thiểu. Tuy nhiên, LIME chọn kernel lân cận và độ thưa
theo heuristic, còn KernelSHAP chọn kernel, hàm mất mát và ràng buộc để nghiệm
thỏa \tn{độ chính xác cục bộ}{local accuracy}, \tn{tính vắng mặt}{missingness} và
\tn{tính nhất quán}{consistency}~\autocite{p10shap}. Sự khác
biệt nằm ở bài toán được giải, không nằm ở việc cả hai đều dùng hồi quy.

\index{KernelSHAP}%
\index{hàm nhân Shapley}%
